\section*{\inTextCircled[10]{9}{15} Great Plaza}\stepcounter{section}\phantomsection\addcontentsline{toc}{section}{Great Plaza}
{\entryfont The streets of Aberdour eventually open onto a vast central plaza before the citadel. Broad terraces, cracked statues, and weathered monuments surround the open space, while several ancient roads converge here from the surrounding districts.}
\begin{DndReadAloud}
	The ruins fall away into an immense stone plaza, its once-polished surface fractured by centuries of neglect.

	At its far end rises the dark bulk of the citadel, overlooking the city from behind monumental walls and sealed gates.

	Along one side of the plaza stands a broad relief carved into the rock. Most of its figures have been worn smooth, but Dwarven warriors, artisans, and scholars can still be distinguished among the faded shapes.

	Beneath them, a second, much smaller line of writing has been added in a noticeably different hand.
\end{DndReadAloud}
\subsection*{Ancient Relief}\stepcounter{subsection}\phantomsection\addcontentsline{toc}{subsection}{Ancient Dwarven Relief}
{\entryfont A broad stone relief dominates one side of the plaza. Its carvings depict the founding and growth of Aberdour: workshops filled with intricate machinery, constructs assisting their creators, and great engines drawing power through the settlement.

Two separate texts accompany the relief. One is written in Dwarvish, the other in Gnomish. Though their wording differs slightly, both describe the same purpose for Aberdour: the pursuit of knowledge, creation, and new means of harnessing power through the combined craft of its inhabitants.

Both passages conclude with a sentence that translates naturally as:}
\begin{DndReadAloud}
	"The mind is the greatest source of power."	
\end{DndReadAloud}
{\noindent\entryfont The apparent agreement hides an important distinction.
\paragraph*{Dwarvish Interpretation} The word translated as mind carries the meaning of spirit, consciousness, or inner will. To the Dwarves, the phrase suggests that true power originates within the living spirit itself.
\paragraph*{Gnomish Interpretation} The corresponding Gnomish word instead refers primarily to knowledge, intellect, or understanding. To the Gnomes, the phrase suggests that power comes from what can be learned, understood, and applied.\\\\
A character familiar with either language can recognize these differing connotations. Magic such as \textbf{Comprehend Languages} conveys the literal meaning of both passages, but the cultural distinction may require closer consideration.}